\chapter{Conventions}
\label{chap:conventions}

\section{Source and tags}

This blueprint follows F.~Han, \emph{An Introduction to Permutation Processes}, version 0.5 (arXiv:2407.09664v1), Chapters 1--5. Page numbers are the printed pages of that version. Every node names the item of the book it comes from, written \han{Theorem 4.4}{74}, and carries one of the following tags.
\begin{itemize}
  \item \fid{as printed} The statement of the book, up to notation.
  \item \fid{adapted} Any change from the printed statement: an added hypothesis or convention, a misprint read as intended, one part of a printed statement, a more general form that the book's proofs use, or the conventions below. In particular, every statement with covering numbers is adapted, because it uses Mathlib's covering numbers, unless it is tagged Mathlib.
  \item \fid{new} Not in the book.
  \item \fid{Mathlib} The statement is a declaration of Mathlib, with Mathlib's conventions. This tag takes precedence over the rule for covering numbers.
\end{itemize}
A node whose content Mathlib already has is marked as such (dark green in the dependency graph), and its text names the Mathlib declarations after ``In Mathlib:''. Its tag still says how the statement relates to the book, so such a node can also be tagged adapted, for example when it states only one part of a printed result.
When a printed statement has to change before it can be formalized, there are two cases. If there is a single fix and it has been checked, the node states the corrected version and is tagged adapted, and a remark says what the book prints and how and why the node differs. If there are several possible fixes, or it is not known whether the statement holds, the node is marked \emph{not ready} (orange in the dependency graph): it quotes the printed statement (``As printed: \ldots''), read with the book's own notation, and a remark says what is missing. Nodes that need material deferred to a later stage are also marked not ready, with a remark.

Some results are only cited in the book. Their nodes name the cited source in a remark. A proof still has to be chosen, from that source or elsewhere, before the node can be assigned.

Labels have the form \texttt{kind:han-c-k} for item $k$ of Chapter $c$ of the book (for example \texttt{thm:han-4-4}), with a suffix for a part of an item. New nodes have descriptive labels.

\section{Conventions for the whole blueprint}

\begin{enumerate}
  \item \textbf{Topology and metric spaces.} All topological, metric and measure-theoretic notions are those of Mathlib, not the definitions printed in \S2.1 of the book.
  \item \textbf{Covering numbers.} We use Mathlib's conventions: covers by closed balls, radius $\epsilon\ge0$ real, values in $\Nat\cup\{\infty\}$, and $\epsilon$-separated meaning pairwise distance $>\epsilon$. The internal covering number $\cN(T,d,\epsilon)$ (centres in $T$) and the external covering number $\cNext(T,d,\epsilon)$ (centres anywhere) are different, and every statement says which one it uses. $\cD(T,d,\epsilon)$ is the packing number.
  \item \textbf{Sub-Gaussian random variables.} Two notions are used: Mathlib's sub-Gaussian moment generating function (\cref{def:subgaussian-mgf}) and the $\psi_2$-norm (\cref{def:han-2-11}). \Cref{lem:subgaussian-bridge} relates them. Chaining (\cref{chap:entropy}) is stated with $\psi$-norms, as in the book.
  \item \textbf{Weak convergence.} Random elements are Borel measurable and weak convergence is Mathlib's convergence in distribution (\cref{def:han-2-2}). In $\linf(T)$ this is \cref{def:linf-weak-conv}; it also asks that the limit law give measure one to $\UC(T,\rho)$. Outer expectations, which the Donsker theorems of \cref{chap:empirical} need, are added later (\cref{def:han-2-8}). An outer probability is the measure of an arbitrary set; Mathlib's measures are defined on all sets.
  \item \textbf{Finite populations.} In \cref{chap:permutation-processes} the probability space is the finite set $\SN$ of permutations of $[N]$ with all subsets measurable and the uniform measure (\cref{def:uniform-permutation}). Every map on it is measurable.
  \item \textbf{Asymptotics.} Limits in \cref{chap:permutation-processes} are along $N\to\infty$ with $n=n_N\to\infty$. The book does not say how the permutations $\pi_N$ for different $N$ are realized together, so its almost-sure statements about them are stated as complete convergence: $\sum_N\Pr(|Y_N|>\epsilon)<\infty$ for every $\epsilon>0$. By the Borel--Cantelli lemma (\cref{lem:borel-cantelli}) this gives $Y_N\to0$ almost surely for every joint realization. Almost-sure statements about a single i.i.d.\ sequence (\cref{chap:empirical}) are unchanged.
  \item \textbf{Entropy integrals.} An entropy integral $\int_0^D\phi\bigl(\cN(T,d,\epsilon)\bigr)\,d\epsilon$ has an integrand with values in $[0,\infty]$ and is a Lebesgue integral of a $[0,\infty]$-valued function (Mathlib's lower Lebesgue integral), so it may be infinite. Covering numbers take values in $\Nat\cup\{\infty\}$, and $\phi(\infty)=\infty$; for example $\psi^{-1}(\infty)=\infty$. The upper limit $D$ is the diameter with values in $[0,\infty]$, so an unbounded set has $D=\infty$. A bound whose right-hand side is infinite holds trivially. A real-valued form of an entropy integral is used only after the integral has been shown to be finite. The logarithm on $[0,\infty]$ takes values in $[-\infty,\infty]$, with $\log0=-\infty$; so integrands such as $\sqrt{\log2\cN}$ are defined separately as $[0,\infty]$-valued functions of $\cN\in\Nat\cup\{\infty\}$, with the value $0$ at $\cN=0$ (an empty set) and $\infty$ at $\cN=\infty$. Expectations of $[0,\infty]$-valued random variables are integrals of the same kind; for a function not known to be measurable, the lower integral is not an outer expectation, so the measurability of a supremum is shown separately (next item).
  \item \textbf{Suprema of processes.} The chaining bounds of \cref{chap:entropy} and their corollaries are stated for processes indexed by a countable pseudo-metric space. Every supremum over the index set, or over pairs of indices, is then a countable supremum of random variables, so it is measurable, with values in $[0,\infty]$. \Cref{lem:chaining-transfer} extends the bounds to a separable pseudo-metric space $(T,d)$ when the paths $t\mapsto X(t)$ are continuous for $d$. The separability of \cref{def:han-2-5} is not used for this, because on its own it does not make the supremum over pairs at distance at most $\delta$ measurable. In \cref{chap:permutation-processes} the probability space is finite and measurability is automatic; there \cref{lem:chaining-transfer} serves to pass to classes that are not countable.
  \item \textbf{Classes of functions in $L^p(Q)$.} A class of functions keeps its functions as indices and evaluates them pointwise; it is not replaced by a set of equivalence classes. For $1\le p\le\infty$, the $L^p(Q)$ distance on a class is the pseudo-metric induced by $f\mapsto[f]\in L^p(Q)$ (\cref{def:lp-class-pseudometric}), and covering and packing numbers of the class are those of \cref{def:han-4-1} for this pseudo-metric. When two such distances are used on one class, for example $L^2(\PN)$ and $L^2(\mathrm P)$ in \cref{chap:permutation-processes}, each is put on its own copy of the index set, tagged by $(p,Q)$. Every $L^p$ in this blueprint has $p\ge1$.
\end{enumerate}
